% Generado por sistema/backend ReporteTex.java -- no editar a mano.
% Datos: Base de datos: 148 importado (titulados) / Base de datos: 53 importado (empleadores)
\subsubsection*{C8. Edad $\times$ área de interés del posgrado}

\textit{Paso 1. Describir el problema.} Analizar el patrón demográfico: establecer si el área en que el titulado quiere hacer un posgrado depende de su edad, para orientar la oferta a cada grupo.

\begin{itemize}[nosep]
  \item Variable independiente $X$: Edad.
  \item Variable dependiente $Y$: Área de interés.
  \item Población: interesados en posgrado ($N = 123$). Criterio ARCU-SUR: 2.4.1. Nivel de significancia $\alpha = 0.05$.
  \item $H_0$: Las variables «Edad» (X) y «Área de interés» (Y) son independientes.
  \item $H_1$: Existe dependencia significativa entre las variables «Edad» (X) y «Área de interés» (Y).
\end{itemize}

\textit{Datos.} Primeros 10 de los 123 titulados que respondieron ambas preguntas (sin datos personales).

\begin{mitabla}{C8: datos de los primeros titulados}
\begin{tabular}{lll}
\toprule
Titulado & $X$: Edad & $Y$: Área de interés \\
\midrule
T1 & 27–34 & Redes y Seguridad \\
T2 & 27–34 & IA y Robótica \\
T3 & 35 o más & IA y Robótica \\
T4 & 26 o menos & Software y Base de Datos \\
T5 & 35 o más & IA y Robótica \\
T6 & 35 o más & IA y Robótica \\
T7 & 27–34 & Software y Base de Datos \\
T8 & 35 o más & Software y Base de Datos \\
T9 & 27–34 & Redes y Seguridad \\
T10 & 27–34 & IA y Robótica \\
\bottomrule
\end{tabular}
\end{mitabla}

\textit{Paso 2. Tabla de contingencia} (frecuencia observada $FO_{ij}$, $N = 123$). Cada celda cuenta a los titulados que respondieron a la vez la fila ($X$) y la columna ($Y$).

\begin{mitabla}{C8: tabla de contingencia, frecuencia observada $FO_{ij}$}
\begin{tabular}{L{3.8cm}C{3.00cm}C{3.00cm}C{3.00cm}|c}
\toprule
$X$: Edad \textbackslash{} $Y$: Área de interés & IA y Robótica & Software y Base de Datos & Redes y Seguridad & Total \\
\midrule
26 o menos & 12 & 18 & 4 & 34 \\
27–34 & 15 & 8 & 15 & 38 \\
35 o más & 32 & 13 & 6 & 51 \\
\midrule
Total & 59 & 39 & 25 & $N = 123$ \\
\bottomrule
\end{tabular}
\end{mitabla}

\textit{Paso 3. Frecuencia esperada.} Es lo que habría en cada celda si $X$ e $Y$ fueran independientes. Por ejemplo, para la celda 26 o menos y IA y Robótica:
\[
FE_{11} = \frac{(34)(59)}{123} = 16.31
\]

\begin{mitabla}{C8: frecuencia esperada $FE_{ij}$}
\begin{tabular}{L{3.8cm}C{3.00cm}C{3.00cm}C{3.00cm}}
\toprule
$X$: Edad \textbackslash{} $Y$: Área de interés & IA y Robótica & Software y Base de Datos & Redes y Seguridad \\
\midrule
26 o menos & $34\cdot59/123 = 16.31$ & $34\cdot39/123 = 10.78$ & $34\cdot25/123 = 6.91$ \\
27–34 & $38\cdot59/123 = 18.23$ & $38\cdot39/123 = 12.05$ & $38\cdot25/123 = 7.72$ \\
35 o más & $51\cdot59/123 = 24.46$ & $51\cdot39/123 = 16.17$ & $51\cdot25/123 = 10.37$ \\
\bottomrule
\end{tabular}
\end{mitabla}

\textit{Paso 4. Estadístico chi-cuadrado.} En cada celda se calcula $\frac{(FO_{ij}-FE_{ij})^2}{FE_{ij}}$ y luego se suman todas.

\begin{mitabla}{C8: cálculo de $X_0^2$ celda por celda}
\begin{tabular}{L{3.2cm}C{3.20cm}C{3.20cm}C{3.20cm}|c}
\toprule
$X$: Edad \textbackslash{} $Y$: Área de interés & IA y Robótica & Software y Base de Datos & Redes y Seguridad & Suma fila \\
\midrule
26 o menos & $\dfrac{(12-16.31)^2}{16.31}$\newline $= 1.1385$ & $\dfrac{(18-10.78)^2}{10.78}$\newline $= 4.8348$ & $\dfrac{(4-6.91)^2}{6.91}$\newline $= 1.2259$ & $7.1991$ \\[6pt]
27–34 & $\dfrac{(15-18.23)^2}{18.23}$\newline $= 0.5715$ & $\dfrac{(8-12.05)^2}{12.05}$\newline $= 1.3605$ & $\dfrac{(15-7.72)^2}{7.72}$\newline $= 6.8552$ & $8.7872$ \\[6pt]
35 o más & $\dfrac{(32-24.46)^2}{24.46}$\newline $= 2.3218$ & $\dfrac{(13-16.17)^2}{16.17}$\newline $= 0.6217$ & $\dfrac{(6-10.37)^2}{10.37}$\newline $= 1.8388$ & $4.7823$ \\[6pt]
\bottomrule
\end{tabular}
\end{mitabla}

\[
X_0^2 = 7.1991 + 8.7872 + 4.7823 = 20.7687
\]

\textit{Paso 5. Estadístico teórico.} Con $f = 3$ filas y $c = 3$ columnas: $gl = (3-1)(3-1) = 4$. De la tabla chi-cuadrada, con $\alpha = 0.05$: $X^2_{4,\,0.05} = 9.488$.

\textit{Paso 6. Relación entre estadísticos y conclusión.}

\begin{center}
\begin{tabular}{l}
$20.7687 \ge 9.488 \rightarrow$ falso\\
Si ($X_0^2 < X^2_{(f-1)(c-1),\,\alpha}$)\\
\quad ``Independencia''\\
Sino\\
\quad \fbox{NO ``Independencia''}\\
FinSi
\end{tabular}
\end{center}

Se rechaza $H_0$ y se acepta $H_1$: «Área de interés» depende de «Edad». Porcentaje de cada fila en la columna ``IA y Robótica'': 26 o menos $12/34 = 35.3\%$; 27–34 $15/38 = 39.5\%$; 35 o más $32/51 = 62.7\%$.

\textit{Recomendación.} El área de interés cambia con la edad: orientar la oferta de posgrado por grupo de edad, con el área que más elige cada grupo (ver la tabla de contingencia).

